\chapter{AGF-ADNet模型设计}

\section{问题定义与总体框架}

面向电源系统多采样率异常检测任务，设原始监测序列记为
\begin{equation}
    \mathbf{D} = [\mathbf{d}_1,\mathbf{d}_2,\ldots,\mathbf{d}_T]^\top \in \mathbb{R}^{T \times N},
\end{equation}
其中，$T$ 表示时间长度，$N$ 表示监测变量数，$\mathbf{d}_t \in \mathbb{R}^{N}$ 表示时刻 $t$ 的多变量观测向量。针对多采样率场景中普遍存在的缺失问题，引入缺失掩码
\begin{equation}
    M_{t,n} =
    \begin{cases}
        1, & \text{第 $t$ 个时刻第 $n$ 个变量缺失},\\
        0, & \text{第 $t$ 个时刻第 $n$ 个变量可观测},
    \end{cases}
\end{equation}
并定义观测掩码为
\begin{equation}
    O_{t,n} = 1 - M_{t,n}.
\end{equation}

本文采用长度为 $L$ 的滑动窗口构造样本，窗口张量表示为
\begin{equation}
    \mathbf{X} \in \mathbb{R}^{B \times L \times N}, \qquad
    \mathbf{M} \in \{0,1\}^{B \times L \times N},
\end{equation}
其中 $B$ 为批大小。在当前实验设置中，窗口长度取 $L=50$，变量总数为 $N=18$。根据各变量观测频率差异，可进一步将变量划分为若干采样率组，并为每个变量赋予采样率类别标识 $r_n$ 与采样步长 $s_n$。

AGF-ADNet 采用插补与重构联合建模的总体思路。多采样率异常检测面临三个彼此耦合的关键问题：缺失位置会削弱时序连续性，变量之间存在非平稳耦合关系，不同采样频率下的周期结构和相位信息又难以通过统一的时域模型充分表达。若仅将缺失值恢复视为一个独立的数据清洗步骤，则插补结果往往只追求点值误差最小，而未能兼顾后续异常判别所需的结构信息；若直接在缺失序列上执行异常检测，则模型又容易将采样率差异和缺失模式混同为异常信号。基于此，AGF-ADNet 将缺失恢复、结构建模和异常判别纳入同一优化框架，使插补结果不仅需要接近真实观测，还需要有利于后续全局重构和异常区分。

从模块构成上看，模型首先利用双专家协同插补模块恢复缺失位置，得到完整窗口；随后，引入采样率感知 Transformer 对完整窗口进行全局重构；最后，结合重构偏差、专家分歧和分布漂移构建异常分数。这样的分层设计具有明确的功能分工：前一层负责恢复完整、稳定且具有结构一致性的输入表示，后一层负责从完整窗口中提取正常工况的全局演化规律。模型的整体映射可概括为
\begin{align}
    \mathbf{X}^{\mathrm{in}} &= \mathbf{X} \odot (1-\widetilde{\mathbf{M}}), \\
    \mathbf{X}^{\mathrm{seed}} &= \mathcal{W}\!\left(\mathbf{X}^{\mathrm{in}}, \widetilde{\mathbf{M}}\right), \\
    \mathbf{A} &= \mathcal{G}\!\left(\mathbf{X}^{\mathrm{seed}}, \widetilde{\mathbf{M}}\right), \\
    \mathbf{X}^{\mathrm{gcn}} &= \mathcal{F}_{\mathrm{gcn}}\!\left(\mathbf{X}^{\mathrm{seed}}, \mathbf{A}\right), \\
    \mathbf{X}^{\mathrm{freq}} &= \mathcal{F}_{\mathrm{freq}}\!\left(\mathbf{X}^{\mathrm{seed}}\right), \\
    \mathbf{X}^{\mathrm{imp}} &= \mathcal{H}\!\left(\mathbf{X}^{\mathrm{seed}}, \widetilde{\mathbf{M}}, \mathbf{X}^{\mathrm{gcn}}, \mathbf{X}^{\mathrm{freq}}, \mathbf{r}\right), \\
    \mathbf{X}^{\mathrm{comp}} &= \mathbf{X}^{\mathrm{in}} + \widetilde{\mathbf{M}} \odot \mathbf{X}^{\mathrm{imp}}, \\
    \hat{\mathbf{X}} &= \mathcal{D}\!\left(\mathbf{X}^{\mathrm{comp}}, \mathbf{O}, \mathbf{r}, \mathbf{s}\right),
\end{align}
其中，$\widetilde{\mathbf{M}}$ 表示模型训练或推理阶段所采用的当前缺失掩码，$\mathbf{A}$ 为自适应学习到的变量图，$\mathbf{X}^{\mathrm{imp}}$ 为融合插补结果，$\mathbf{X}^{\mathrm{comp}}$ 为完整窗口，$\hat{\mathbf{X}}$ 为最终重构输出。

为了便于后续分析，表~\ref{tab:agf_adnet_modules} 给出了 AGF-ADNet 主要模块的功能划分及输入输出形式。

\begin{table}[htbp]
    \centering
    \caption{AGF-ADNet 主要模块及输入输出}
    \label{tab:agf_adnet_modules}
    \begin{tabular}{p{2.4cm} p{4.8cm} p{3.0cm} p{3.0cm}}
        \hline
        模块 & 主要功能 & 输入 & 输出 \\
        \hline
        数据预处理 & 完成缺失掩码生成、观测值标准化、窗口化与采样率元数据推断 & 原始序列 $\mathbf{D}$ & $\mathbf{X}$、$\mathbf{M}$、$\mathbf{r}$、$\mathbf{s}$ \\
        双向预插补 & 为后续学习型插补提供连续初值 & $\mathbf{X}^{\mathrm{in}}$、$\widetilde{\mathbf{M}}$ & $\mathbf{X}^{\mathrm{seed}}$ \\
        自适应图学习 & 依据窗口统计量学习样本相关的变量邻接结构 & $\mathbf{X}^{\mathrm{seed}}$、$\widetilde{\mathbf{M}}$ & $\mathbf{A}$ \\
        图重构专家 & 利用变量图执行跨变量信息传播与重构 & $\mathbf{X}^{\mathrm{seed}}$、$\mathbf{A}$ & $\mathbf{X}^{\mathrm{gcn}}$ \\
        频域重构专家 & 利用频谱增强提取周期性与谐波信息 & $\mathbf{X}^{\mathrm{seed}}$ & $\mathbf{X}^{\mathrm{freq}}$ \\
        采样率感知门控融合 & 动态分配双专家权重并形成缺失值估计 & $\mathbf{X}^{\mathrm{seed}}$、$\widetilde{\mathbf{M}}$、$\mathbf{X}^{\mathrm{gcn}}$、$\mathbf{X}^{\mathrm{freq}}$、$\mathbf{r}$ & $\mathbf{X}^{\mathrm{imp}}$、$\mathbf{X}^{\mathrm{comp}}$ \\
        采样率感知重构检测器 & 建模完整窗口的全局时序模式并输出重构序列 & $\mathbf{X}^{\mathrm{comp}}$、$\mathbf{O}$、$\mathbf{r}$、$\mathbf{s}$ & $\hat{\mathbf{X}}$ \\
        异常评分与判别 & 融合多源统计量并完成阈值判别 & $\hat{\mathbf{X}}$、$\mathbf{X}^{\mathrm{comp}}$、$\mathbf{X}^{\mathrm{gcn}}$、$\mathbf{X}^{\mathrm{freq}}$ & 异常分数与预测结果 \\
        \hline
    \end{tabular}
\end{table}

\section{数据处理与样本构造}

\subsection{缺失掩码生成与观测值标准化}

多采样率工业数据在统一时间轴对齐后，不同变量会在非采样时刻形成大量空缺项。本文以缺失掩码 $\mathbf{M}$ 显式记录该信息，并仅基于可观测位置估计标准化参数。对于第 $n$ 个变量，其均值与标准差定义为
\begin{align}
    \mu_n &= \frac{\sum_{t=1}^{T} O_{t,n} D_{t,n}}{\sum_{t=1}^{T} O_{t,n} + \varepsilon}, \\
    \sigma_n &= \sqrt{\frac{\sum_{t=1}^{T} O_{t,n}(D_{t,n}-\mu_n)^2}{\sum_{t=1}^{T} O_{t,n} + \varepsilon} + \varepsilon},
\end{align}
其中 $\varepsilon$ 为防止分母为零的极小常数。随后，将缺失位置临时以变量均值替代，得到
\begin{equation}
    \bar{D}_{t,n} = O_{t,n} D_{t,n} + M_{t,n} \mu_n,
\end{equation}
并执行标准化变换
\begin{equation}
    \tilde{D}_{t,n} = \frac{\bar{D}_{t,n}-\mu_n}{\sigma_n}.
\end{equation}

这一处理方式具有两点优势：其一，标准化参数仅由真实观测值确定，避免缺失模式对统计量造成偏移；其二，缺失位置在标准化后恰好变为零，从而可与显式掩码协同输入网络。此外，采用观测值均值进行占位还可以保证不同变量在输入端具有统一的数值尺度，使后续图学习、频域变换和门控融合不会因为缺失占位值过大或过小而产生额外偏差。对于图结构学习而言，零值占位配合显式掩码能够将数值信息与可观测性信息分离；对于频域分支而言，经过标准化后的输入具有更稳定的频谱幅值范围，有利于注意力权重与频谱增量的联合学习。

\subsection{滑动窗口构造}

为兼顾局部动态建模与序列级异常评分，本文采用两类窗口化策略。

\begin{enumerate}
    \item 对于插补训练与全序列补全过程，采用前端填充滑动窗口。若某一时刻历史长度不足 $L$，则以前序首样本进行补齐，使每一个时间索引都对应一个长度为 $L$ 的输入窗口。
    \item 对于异常评分阶段，采用标准滑动窗口，即仅在历史长度达到 $L$ 后开始取窗，以避免前端填充对统计分数带来额外偏差。
\end{enumerate}

以前端填充窗口为例，第 $t$ 个时间点对应的窗口样本可表示为
\begin{equation}
    \mathbf{X}^{(t)} =
    \begin{cases}
        [\underbrace{\tilde{\mathbf{d}}_1,\ldots,\tilde{\mathbf{d}}_1}_{L-t},\tilde{\mathbf{d}}_1,\ldots,\tilde{\mathbf{d}}_{t}], & 1 \leq t < L, \\
        [\tilde{\mathbf{d}}_{t-L+1},\ldots,\tilde{\mathbf{d}}_{t}], & L \leq t \leq T,
    \end{cases}
\end{equation}
相应掩码窗口按完全相同的方式构造。这样既保证了补全过程与原始时间轴严格对齐，又为模型提供了统一长度的输入张量。前端填充窗口主要服务于补全输出与原序列的逐点对齐，因为模型最终需要为序列中的每一个时间位置提供重构或补全结果；标准窗口则主要服务于异常评分阶段，因为该阶段更关注具有完整历史支撑的统计特征。二者分别面向补全过程与判别过程的不同需求，有助于兼顾时间对齐精度和评价稳定性。

\subsection{采样率元数据推断}

对于多采样率过程，变量的观测比例本身包含采样步长信息。设第 $n$ 个变量的观测比例为
\begin{equation}
    \rho_n = \frac{1}{T}\sum_{t=1}^{T} O_{t,n},
\end{equation}
则其采样步长可近似写为
\begin{equation}
    s_n = \operatorname{round}\!\left(\frac{1}{\max(\rho_n,\varepsilon)}\right).
\end{equation}
根据不同步长的取值，可进一步映射为离散采样率类别标识 $r_n$。在当前实验数据中，18 个变量可划分为三个采样率组，其步长分别为 $1$、$3$ 和 $5$，每组包含 6 个变量。该元数据在后续门控融合与 Transformer 编码中共同发挥作用。

\subsection{自监督再遮挡机制}

原始结构性缺失位置不具备真实标签，因此插补训练不能直接以其作为监督目标。为此，本文在训练阶段从原本可观测的位置中随机抽取一部分进行再遮挡，记为 $\mathbf{R}$。于是模型实际看到的缺失掩码为
\begin{equation}
    \widetilde{\mathbf{M}} = \max(\mathbf{M}, \mathbf{R}),
\end{equation}
其中 $\mathbf{R}_{t,n}=1$ 表示该位置在训练阶段被额外遮挡。进一步得到模型输入
\begin{equation}
    \mathbf{X}^{\mathrm{in}} = \mathbf{X} \odot (1-\widetilde{\mathbf{M}}).
\end{equation}
这种自监督构造使模型可以在具有真实值的位置上学习缺失恢复能力，同时保持结构性缺失场景下的训练一致性。进一步地，再遮挡机制还迫使模型在不同缺失模式下学习稳定的恢复映射，从而降低模型对固定缺失分布的过拟合风险。由于训练阶段的人为遮挡仅发生在原本可观测的位置，插补误差始终可以与真实值对齐计算，因此训练目标具有明确的统计意义。

\section{双专家协同插补模块}

\subsection{双向预插补初始化}

直接将零填充序列送入学习型插补器，容易造成真实零值与占位零值的混淆。为减小这一不利影响，AGF-ADNet 在双专家建模之前加入双向预插补初始化。预插补的目标不是直接得到最终估计，而是先恢复一个具有基本时间连续性的初始序列，使后续学习模块在进入图卷积传播和频域分析之前，能够获得更稳定的数值基础。设当前窗口中第 $n$ 个变量在时刻 $t$ 的前向最近观测位置和后向最近观测位置分别为
\begin{align}
    \tau_{t,n}^{\mathrm{f}} &= \max\{k \mid k \leq t,\ \widetilde{M}_{k,n}=0\}, \\
    \tau_{t,n}^{\mathrm{b}} &= \min\{k \mid k \geq t,\ \widetilde{M}_{k,n}=0\}.
\end{align}
据此可构造前向与后向填充值
\begin{align}
    x_{t,n}^{\mathrm{f}} &= X_{\tau_{t,n}^{\mathrm{f}},n}, \\
    x_{t,n}^{\mathrm{b}} &= X_{\tau_{t,n}^{\mathrm{b}},n}.
\end{align}
双向预插补结果定义为
\begin{equation}
    X_{t,n}^{\mathrm{seed}} =
    \begin{cases}
        X_{t,n}, & \widetilde{M}_{t,n}=0, \\
        \dfrac{x_{t,n}^{\mathrm{f}} + x_{t,n}^{\mathrm{b}}}{2}, & x_{t,n}^{\mathrm{f}},x_{t,n}^{\mathrm{b}} \text{ 均存在}, \\
        x_{t,n}^{\mathrm{f}}, & x_{t,n}^{\mathrm{f}} \text{ 存在而 } x_{t,n}^{\mathrm{b}} \text{ 不存在}, \\
        x_{t,n}^{\mathrm{b}}, & x_{t,n}^{\mathrm{b}} \text{ 存在而 } x_{t,n}^{\mathrm{f}} \text{ 不存在}, \\
        0, & \text{前后均不存在有效观测}.
    \end{cases}
\end{equation}

双向预插补并不追求最终最优插补，而是为后续图专家与频域专家提供更平滑、更接近真实演化轨迹的初始化序列。之所以同时采用前向与后向两类信息，是因为单向保持策略会对长缺失段形成系统性偏置。仅使用前向最近观测值时，序列会在缺失区间内呈现明显阶梯状保持；仅使用后向最近观测值时，则会引入相反方向的偏差。双向预插补通过综合缺失区间两侧的信息，能够在不引入额外参数的前提下缓和上述偏差，并降低后续学习模块将缺失边界误识别为异常结构的风险。

\subsection{自适应图学习机制}

工业过程中的多变量之间通常存在复杂的耦合与传递关系。为在不依赖先验物理拓扑的条件下刻画这种结构依赖，本文设计自适应图学习模块，为每个输入窗口动态生成变量邻接矩阵。该模块的核心出发点是，工业变量之间的有效关联并非全局固定，而会随着当前运行状态、缺失比例以及局部统计特征的变化而变化。因此，图结构应当由当前窗口内的状态统计驱动，而不是由预先设定的固定拓扑完全决定。

对于第 $b$ 个窗口样本中的第 $n$ 个变量，首先提取四类动态统计量：
\begin{align}
    \mu_{b,n} &= \frac{\sum_{t=1}^{L} O_{b,t,n}^{\prime} X_{b,t,n}^{\mathrm{seed}}}{\sum_{t=1}^{L} O_{b,t,n}^{\prime} + \varepsilon}, \\
    \sigma_{b,n} &= \sqrt{\frac{\sum_{t=1}^{L} O_{b,t,n}^{\prime} \left(X_{b,t,n}^{\mathrm{seed}}-\mu_{b,n}\right)^2}{\sum_{t=1}^{L} O_{b,t,n}^{\prime} + \varepsilon} + \varepsilon}, \\
    \ell_{b,n} &= X_{b,\tau_{b,n}^{\ast},n}^{\mathrm{seed}}, \qquad \tau_{b,n}^{\ast} = \arg\max_t \left(t \cdot O_{b,t,n}^{\prime}\right), \\
    r_{b,n}^{\mathrm{miss}} &= \frac{1}{L}\sum_{t=1}^{L} \widetilde{M}_{b,t,n},
\end{align}
其中 $O_{b,t,n}^{\prime}=1-\widetilde{M}_{b,t,n}$ 表示当前模型视角下的观测指示。上述四类统计量分别对应不同的信息来源：均值与标准差反映窗口内的总体分布状态，最近一次观测值强调当前时刻附近的运行水平，缺失比例反映该变量在当前窗口中的信息可靠性。于是可得节点的动态上下文向量
\begin{equation}
    \mathbf{z}_{b,n}^{\mathrm{dyn}} =
    [\mu_{b,n},\sigma_{b,n},\ell_{b,n},r_{b,n}^{\mathrm{miss}}].
\end{equation}

进一步地，为每个变量引入静态上下文向量 $\mathbf{c}_n$ 与潜在坐标向量 $\mathbf{u}_n$，分别描述变量的稳定身份信息与潜在几何关系。静态上下文保证不同变量在长期意义上具有可区分的身份表征，潜在坐标则为变量间的先验接近程度提供连续参数化形式。将动态上下文与静态上下文拼接后，经多层感知机编码为节点表示
\begin{equation}
    \mathbf{h}_{b,n} = \phi_{\mathrm{node}}\!\left([\mathbf{z}_{b,n}^{\mathrm{dyn}},\mathbf{c}_n]\right).
\end{equation}

为了估计变量 $i$ 与变量 $j$ 之间的关联强度，构造如下配对特征：
\begin{equation}
    \mathbf{p}_{b,i,j} =
    [\mathbf{h}_{b,i},\mathbf{h}_{b,j},|\mathbf{h}_{b,i}-\mathbf{h}_{b,j}|,\mathbf{h}_{b,i}\odot\mathbf{h}_{b,j}],
\end{equation}
并通过边权映射得到边对数值
\begin{equation}
    e_{b,i,j} = \phi_{\mathrm{edge}}(\mathbf{p}_{b,i,j}).
\end{equation}

与此同时，依据潜在坐标构造距离先验
\begin{equation}
    P_{i,j} =
    \begin{cases}
        \dfrac{1}{1+\|\mathbf{u}_i-\mathbf{u}_j\|_2}, & i \neq j, \\
        0, & i=j,
    \end{cases}
\end{equation}
并利用归一化节点表示计算余弦相似度
\begin{equation}
    S_{b,i,j} = \frac{\mathbf{h}_{b,i}^{\top}\mathbf{h}_{b,j}}{\|\mathbf{h}_{b,i}\|_2 \|\mathbf{h}_{b,j}\|_2}.
\end{equation}

综合以上三类信息后，邻接矩阵的未归一化对数值写为
\begin{equation}
    L_{b,i,j} = e_{b,i,j} + \alpha P_{i,j} + \beta S_{b,i,j} + \gamma \delta_{i,j},
\end{equation}
其中，$\alpha$、$\beta$、$\gamma$ 为可学习系数，$\delta_{i,j}$ 为 Kronecker 符号。随后引入温度参数 $\tau>0$ 执行 softmax 归一化，并通过对称化与行归一化得到最终邻接矩阵：
\begin{align}
    \bar{\mathbf{A}}_b &= \operatorname{softmax}\!\left(\frac{\mathbf{L}_b}{\tau}\right), \\
    \tilde{\mathbf{A}}_b &= \frac{1}{2}\left(\bar{\mathbf{A}}_b + \bar{\mathbf{A}}_b^\top\right), \\
    A_{b,i,j} &= \frac{\tilde{A}_{b,i,j}}{\sum_{j=1}^{N} \tilde{A}_{b,i,j} + \varepsilon}.
\end{align}

由此获得的 $\mathbf{A}_b \in \mathbb{R}^{N \times N}$ 为样本相关动态图，能够随窗口状态变化而自适应调整。该图同时包含三类信息来源，即窗口内统计关系、潜在几何先验以及表示相似性，因此可以在缺乏明确物理拓扑先验的条件下形成相对稳定的关系矩阵。对角平衡和温度缩放共同保证了图结构既不过度依赖自连接，也不会在变量之间产生过于尖锐或过于均匀的分配。

\subsection{图重构专家}

图重构专家以学习得到的变量图为基础，通过图卷积实现跨变量信息传播。设第 $b$ 个样本在时刻 $t$ 的输入向量为 $\mathbf{x}_{b,t}^{\mathrm{seed}} \in \mathbb{R}^{N}$，其图卷积更新可写为
\begin{equation}
    \mathbf{H}_{b,t}^{(l+1)} = \mathbf{A}_b \mathbf{H}_{b,t}^{(l)} \mathbf{W}^{(l)},
\end{equation}
其中 $\mathbf{H}_{b,t}^{(0)} = \mathbf{x}_{b,t}^{\mathrm{seed}}$，$\mathbf{W}^{(l)}$ 为第 $l$ 层线性映射参数。AGF-ADNet 采用三层图卷积结构：前两层之后接非线性激活、层归一化与随机失活操作，最后一层回到标量空间，得到图专家输出
\begin{equation}
    \mathbf{X}^{\mathrm{gcn}} \in \mathbb{R}^{B \times L \times N}.
\end{equation}

该分支在整个插补框架中的作用主要体现在两个方面。第一，图卷积在同一时间步沿变量维传播信息，因此能够直接利用变量之间的同步耦合关系恢复缺失项，这对受控制量、反馈量和关联测点之间存在强联动的工业系统尤为重要。第二，图结构由整个窗口的统计状态决定，而传播过程在窗口内各时间步共享同一张图，因此该分支能够在局部数值传播与整窗状态感知之间建立联系。三层图卷积结构既保留了邻近相关变量的直接贡献，也允许更高阶关系在网络中逐步传播，从而提升对复杂变量关联的表达能力。

\subsection{频域重构专家}

与图重构专家强调跨变量依赖不同，频域重构专家主要关注单变量序列中的周期结构、谐波特征和振荡模式。对于任意窗口样本，将输入张量沿时间维进行实数快速傅里叶变换，可得
\begin{equation}
    \mathbf{F}_{b,n} = \operatorname{rFFT}\!\left(\mathbf{X}_{b,:,n}^{\mathrm{seed}}\right)
    = \mathbf{R}_{b,n} + j \mathbf{I}_{b,n},
\end{equation}
其中 $\mathbf{R}_{b,n}$ 和 $\mathbf{I}_{b,n}$ 分别表示实部与虚部。

将实部与虚部拼接为频谱表征向量
\begin{equation}
    \mathbf{q}_{b,n} = [\mathbf{R}_{b,n},\mathbf{I}_{b,n}],
\end{equation}
并构建频率注意力
\begin{equation}
    \boldsymbol{\alpha}_{b,n} = \sigma\!\left(\phi_{\mathrm{att}}(\mathbf{q}_{b,n})\right),
\end{equation}
其中 $\sigma(\cdot)$ 表示 Sigmoid 函数。进一步，频谱增强网络输出实部和虚部的增量项
\begin{equation}
    [\Delta \mathbf{R}_{b,n}, \Delta \mathbf{I}_{b,n}] = \phi_{\mathrm{enh}}(\mathbf{q}_{b,n}).
\end{equation}
于是增强后的复频谱写为
\begin{equation}
    \tilde{\mathbf{F}}_{b,n}
    =
    \mathbf{F}_{b,n}
    +
    \boldsymbol{\alpha}_{b,n} \odot
    \left(
        \Delta \mathbf{R}_{b,n} + j \Delta \mathbf{I}_{b,n}
    \right).
\end{equation}

通过逆傅里叶变换可恢复频域增强后的时域序列
\begin{equation}
    \mathbf{X}_{b,:,n}^{\prime} = \operatorname{irFFT}(\tilde{\mathbf{F}}_{b,n}),
\end{equation}
随后再经共享前馈映射头进行残差修正，得到频域专家输出
\begin{equation}
    \mathbf{X}^{\mathrm{freq}} = \mathbf{X}^{\prime} + \phi_{\mathrm{out}}(\mathbf{X}^{\prime}).
\end{equation}

该分支能够显式强化频谱特征，对电压、电流等具有周期性波动规律的信号具有较强的适应能力。其设计关键在于直接对频谱实部和虚部进行联合建模，而不是仅利用幅值谱。这样既能够保留周期强度信息，也能保留相位相关信息，从而更完整地描述周期模式。采用残差形式修正频谱而不是完全重构频谱，有助于维持原始频谱的稳定结构，避免频域分支在训练初期产生过度振荡。

\subsection{采样率感知门控融合}

由于图专家与频域专家擅长刻画的信息类型不同，简单平均往往难以适应变量异质性和时变工况。为此，AGF-ADNet 引入采样率感知门控融合模块，对每个时间步、每个变量独立生成双专家权重。该设计的核心思想是，将专家选择问题转化为条件加权问题，使模型能够根据变量采样频率、当前位置的可观测性以及两类专家相对预插补结果的偏移程度，自适应地决定更可信的恢复路径。

首先，为每个变量构造采样率嵌入向量 $\mathbf{e}_n^{r}$。然后，对位置 $(b,t,n)$ 形成门控输入特征
\begin{equation}
    \mathbf{g}_{b,t,n}^{\mathrm{in}} =
    \left[
        X_{b,t,n}^{\mathrm{seed}},
        \widetilde{M}_{b,t,n},
        X_{b,t,n}^{\mathrm{gcn}},
        X_{b,t,n}^{\mathrm{freq}},
        X_{b,t,n}^{\mathrm{gcn}} - X_{b,t,n}^{\mathrm{seed}},
        X_{b,t,n}^{\mathrm{freq}} - X_{b,t,n}^{\mathrm{seed}},
        \mathbf{e}_n^{r}
    \right].
\end{equation}
经多层感知机后得到双专家门控系数
\begin{equation}
    [\omega_{b,t,n}^{\mathrm{gcn}},\omega_{b,t,n}^{\mathrm{freq}}]
    =
    \operatorname{softmax}\!\left(\phi_{\mathrm{gate}}(\mathbf{g}_{b,t,n}^{\mathrm{in}})\right),
\end{equation}
并满足
\begin{equation}
    \omega_{b,t,n}^{\mathrm{gcn}} + \omega_{b,t,n}^{\mathrm{freq}} = 1.
\end{equation}

融合插补结果定义为
\begin{equation}
    X_{b,t,n}^{\mathrm{imp}}
    =
    \omega_{b,t,n}^{\mathrm{gcn}} X_{b,t,n}^{\mathrm{gcn}}
    +
    \omega_{b,t,n}^{\mathrm{freq}} X_{b,t,n}^{\mathrm{freq}}.
\end{equation}
最终，完整窗口由观测值与插补值共同构成：
\begin{equation}
    \mathbf{X}^{\mathrm{comp}}
    =
    \mathbf{X}^{\mathrm{in}}
    +
    \widetilde{\mathbf{M}} \odot \mathbf{X}^{\mathrm{imp}}.
\end{equation}

门控输入中同时包含掩码信息、双专家输出以及相对于预插补序列的偏移量，原因在于这些量分别描述了当前位置的观测状态、两类恢复机理给出的估计结果以及相对于基础连续性假设的修正幅度。采样率嵌入进一步使门控网络能够感知不同频率变量在信息可靠性和动态特征上的差异。最终，模型并非对整个窗口进行无差别替换，而是只在当前缺失位置写入融合估计，在可观测位置严格保留已知信息，从而增强了插补结果的物理一致性，并减少对原始观测的二次扰动。

\section{采样率感知全局重构检测器}

\subsection{设计思想}

缺失值恢复之后，仍需要从完整窗口中判断样本是否偏离正常运行规律。为此，AGF-ADNet 在插补模块之后引入采样率感知 Transformer 重构检测器。该检测器不再直接针对原始缺失输入工作，而是面向完整窗口 $\mathbf{X}^{\mathrm{comp}}$ 学习正常模式下的全局演化规律，再以重构偏差刻画异常程度。之所以将异常检测置于插补之后，是因为多采样率序列在统一时间轴上存在结构性空缺，若直接对不完整输入建模，检测器往往需要同时学习缺失模式、采样频率差异和正常动态规律，任务耦合过强。通过先获得较为完整的窗口表示，检测器可以更集中地刻画系统在正常工况下的全局时序结构。

\subsection{变量级编码与采样相位建模}

为同时体现变量身份、多采样率属性和周期相位位置，检测器对每个变量在每个时间步构造联合表征。设第 $n$ 个变量的变量嵌入和采样率嵌入分别为 $\mathbf{e}_n^{v}$ 与 $\mathbf{e}_n^{r}$。进一步，根据采样步长 $s_n$ 构造相位标量
\begin{equation}
    \pi_{t,n} = \frac{\operatorname{mod}(t,s_n)}{s_n},
\end{equation}
并通过非线性映射得到相位嵌入
\begin{equation}
    \mathbf{e}_{t,n}^{p} = \phi_{\mathrm{phase}}(\pi_{t,n}).
\end{equation}

于是，位置 $(b,t,n)$ 的变量级输入可写为
\begin{equation}
    \mathbf{u}_{b,t,n}
    =
    \left[
        X_{b,t,n}^{\mathrm{comp}},
        O_{b,t,n},
        \mathbf{e}_n^{v},
        \mathbf{e}_n^{r},
        \mathbf{e}_{t,n}^{p}
    \right].
\end{equation}
其中，$O_{b,t,n}$ 表示原始结构性观测指示，而非训练时额外再遮挡后的掩码。这样可使检测器明确区分真实缺测与训练期自监督遮挡两类情况。变量级输入经线性投影后得到隐藏表示
\begin{equation}
    \mathbf{z}_{b,t,n} = \phi_{\mathrm{var}}(\mathbf{u}_{b,t,n}).
\end{equation}

这一变量级编码过程将数值信息与元数据显式融合。变量嵌入用于区分不同测点的语义身份，采样率嵌入用于描述不同变量的采样层级，相位嵌入则用于反映当前时间位置在该变量采样周期中的相对位置。三者共同作用后，检测器可以在统一时间轴上建立更符合多采样率特性的内部表示。

\subsection{时间 Token 构造与 Transformer 重构}

在时间维度上，将原始完整向量、观测指示向量以及全部变量级隐藏表示进行拼接，形成时刻 $t$ 的时间 Token：
\begin{equation}
    \mathbf{z}_{b,t}^{\mathrm{tok}}
    =
    \phi_{\mathrm{tok}}
    \left(
        \left[
            \mathbf{X}_{b,t,:}^{\mathrm{comp}},
            \mathbf{O}_{b,t,:},
            \operatorname{vec}(\mathbf{z}_{b,t,1},\ldots,\mathbf{z}_{b,t,N})
        \right]
    \right).
\end{equation}
为保留时序位置信息，再加入正弦位置编码：
\begin{equation}
    \mathbf{h}_{b,t}^{(0)} = \mathbf{z}_{b,t}^{\mathrm{tok}} + \mathbf{p}_t.
\end{equation}

随后，利用多层 Transformer 编码器学习全局时间依赖
\begin{equation}
    \mathbf{h}_{b}^{(l+1)} = \operatorname{Transformer}\!\left(\mathbf{h}_{b}^{(l)}\right), \qquad l=0,1,\ldots,L_{\mathrm{enc}}-1.
\end{equation}
经重构头输出增量项 $\Delta \mathbf{X}_{b,t}$ 后，最终重构结果采用残差形式给出：
\begin{equation}
    \hat{\mathbf{X}}_{b,t} = \mathbf{X}_{b,t}^{\mathrm{comp}} + \Delta \mathbf{X}_{b,t}.
\end{equation}

在当前实验设置中，检测器采用 3 层 Transformer 编码器，时间 Token 维度为 128，多头注意力头数为 4。Transformer 的优势在于可以直接在整个窗口范围内建模任意两个时间位置之间的依赖关系，因此不仅能够感知局部突变，也能够刻画慢变漂移、周期累积误差以及跨阶段状态转移。采用残差重构头而不是直接输出完整重构结果，有利于让网络重点学习相对于输入完整窗口的修正量，从而提高训练稳定性。

\section{训练目标与优化机制}

\subsection{插补监督目标}

利用训练阶段构造的再遮挡掩码 $\mathbf{R}$，可分别定义融合插补损失、图专家损失与频域专家损失：
\begin{align}
    \mathcal{L}_{\mathrm{fusion}}
    &=
    \frac{\|(\mathbf{X}^{\mathrm{imp}}-\mathbf{X}) \odot \mathbf{R}\|_F^2}{\|\mathbf{R}\|_1 + \varepsilon}, \\
    \mathcal{L}_{\mathrm{gcn}}
    &=
    \frac{\|(\mathbf{X}^{\mathrm{gcn}}-\mathbf{X}) \odot \mathbf{R}\|_F^2}{\|\mathbf{R}\|_1 + \varepsilon}, \\
    \mathcal{L}_{\mathrm{freq}}
    &=
    \frac{\|(\mathbf{X}^{\mathrm{freq}}-\mathbf{X}) \odot \mathbf{R}\|_F^2}{\|\mathbf{R}\|_1 + \varepsilon}.
\end{align}

上述设计确保双专家与融合器均在具备真实参考值的位置上接受监督，从而避免将结构性缺失位置误用为训练标签。需要指出的是，融合插补损失直接约束最终恢复结果，单专家损失则分别约束图专家和频域专家的独立恢复能力。这样的联合监督有助于防止门控网络仅依赖某一路专家，也有助于提升各专家分支在不同工况下的可分辨性。

\subsection{图结构正则项}

自适应图学习虽然具有较强表达能力，但若缺乏适当约束，容易出现过强自环或过度平滑现象。为此，对邻接矩阵对角项施加平衡约束
\begin{equation}
    \mathcal{L}_{\mathrm{graph}}
    =
    \frac{1}{BN}
    \sum_{b=1}^{B}\sum_{n=1}^{N}
    \left(
        A_{b,n,n} - \rho
    \right)^2,
\end{equation}
其中 $\rho$ 表示期望的对角占比。在当前设置中，$\rho=0.25$。该约束使模型在保留节点自信息与利用邻接传播之间取得平衡。

\subsection{重构检测目标}

检测器的训练目标是在正常样本上学习完整窗口的稳定重构模式。为避免对结构性缺失位置施加伪监督，仅在原始可观测位置上计算重构损失：
\begin{equation}
    \mathcal{L}_{\mathrm{det}}
    =
    \frac{\|(\hat{\mathbf{X}}-\mathbf{X}^{\mathrm{comp}}) \odot \mathbf{O}\|_F^2}{\|\mathbf{O}\|_1 + \varepsilon}.
\end{equation}

这里采用完整窗口 $\mathbf{X}^{\mathrm{comp}}$ 作为检测器的自监督目标，其目标是在正常样本分布上建立完整时序结构的稳定重构映射，而非仅进行局部点值回归。由于检测损失仅在原始可观测位置上计算，检测器不会被迫对原本无真值的位置给出伪精确监督，而是通过可观测部分学习正常模式的整体一致性。

\subsection{总损失函数与优化}

综合以上各项，AGF-ADNet 的总损失写为
\begin{equation}
    \mathcal{L}
    =
    \lambda_1 \mathcal{L}_{\mathrm{fusion}}
    +
    \lambda_2 \mathcal{L}_{\mathrm{gcn}}
    +
    \lambda_3 \mathcal{L}_{\mathrm{freq}}
    +
    \lambda_4 \mathcal{L}_{\mathrm{det}}
    +
    \lambda_5 \mathcal{L}_{\mathrm{graph}}.
\end{equation}
在当前实验配置中，权重取值分别为
\begin{equation}
    \lambda_1 = 1.0, \quad
    \lambda_2 = 0.4, \quad
    \lambda_3 = 0.4, \quad
    \lambda_4 = 0.5, \quad
    \lambda_5 = 0.1.
\end{equation}

模型整体采用 AdamW 优化器进行端到端联合训练，并结合梯度裁剪抑制训练过程中的梯度爆炸风险。由于插补器与检测器共享前向链路，检测器对正常模式的建模需求可以反向促进插补结果朝向更利于异常区分的方向优化。换言之，插补模块不仅需要在再遮挡位置恢复真实值，还需要生成能够被全局重构器稳定建模的完整窗口表示；检测模块则在此基础上学习正常工况的时序约束。两部分通过联合优化形成相互约束关系。

\section{异常评分与判别方法}

\subsection{多源异常分数构造}

AGF-ADNet 不将单一重构误差直接作为异常依据，而是联合考虑三类互补信号。这样设计的原因在于，不同异常类型在统计层面呈现出的特征并不相同。某些异常会直接导致全局重构困难，某些异常则会表现为双专家响应不一致，还有一些异常主要体现为窗口统计分布的持续偏移。因此，融合多源分数能够提高模型对不同异常形态的覆盖能力。

\subsubsection{检测重构分数}

首先，在原始可观测位置上统计检测器的平均平方误差，定义为
\begin{equation}
    s_b^{\mathrm{det}}
    =
    \frac{\|(\hat{\mathbf{X}}_b-\mathbf{X}_b^{\mathrm{comp}})\odot \mathbf{O}_b\|_F^2}{\|\mathbf{O}_b\|_1 + \varepsilon}.
\end{equation}
该分数反映窗口整体动态模式与正常重构规律之间的偏离程度，是最终异常评分的主体部分。

\subsubsection{双专家分歧分数}

其次，利用图专家与频域专家的输出差异刻画结构性不确定性：
\begin{equation}
    s_b^{\mathrm{gap}}
    =
    \frac{1}{LN}
    \sum_{t=1}^{L}\sum_{n=1}^{N}
    \left|
        X_{b,t,n}^{\mathrm{gcn}} - X_{b,t,n}^{\mathrm{freq}}
    \right|.
\end{equation}
若某一窗口内两类机理给出的估计显著不一致，往往说明该窗口已偏离正常数据分布，或者存在更复杂的异常耦合模式。

\subsubsection{分布漂移分数}

为进一步增强对缓慢漂移和全局统计变化的感知能力，本文从补全窗口中提取统计特征向量。设当前实验中的三个采样率组分别为
\begin{equation}
    \mathcal{G}_1=\{1,\ldots,6\},\quad
    \mathcal{G}_2=\{7,\ldots,12\},\quad
    \mathcal{G}_3=\{13,\ldots,18\},
\end{equation}
并引入全局变量集合 $\mathcal{G}_4=\{1,\ldots,18\}$。对于每个集合 $\mathcal{G}_k$，分别计算窗口内的时间均值向量和标准差向量；此外，再补充平均绝对一阶差分以及前 5 个非直流频点的平均幅值特征。最终将上述统计量拼接成特征向量 $\mathbf{f}_b$，并基于训练集统计量进行标准化：
\begin{equation}
    \mathbf{z}_b^{\mathrm{shift}} = \frac{\mathbf{f}_b-\boldsymbol{\mu}_f}{\boldsymbol{\sigma}_f}.
\end{equation}
据此定义分布漂移分数
\begin{equation}
    s_b^{\mathrm{shift}}
    =
    \sqrt{
        \frac{1}{d_f}
        \sum_{i=1}^{d_f}
        \left(z_{b,i}^{\mathrm{shift}}\right)^2
    },
\end{equation}
其中 $d_f$ 表示统计特征维度。该分数对缓变型异常与整体工况漂移具有较好敏感性。

\subsection{分数归一化与融合}

由于三类分数量纲不同，首先利用训练集分布进行标准化：
\begin{equation}
    \bar{s}^{(\cdot)} = \frac{s^{(\cdot)}-\mu^{(\cdot)}_{\mathrm{tr}}}{\sigma^{(\cdot)}_{\mathrm{tr}} + \varepsilon}.
\end{equation}
最终异常分数定义为
\begin{equation}
    s_b
    =
    \left|\bar{s}_b^{\mathrm{det}}\right|
    +
    \eta_1 \bar{s}_b^{\mathrm{gap}}
    +
    \eta_2 \bar{s}_b^{\mathrm{shift}},
\end{equation}
其中，在当前配置中 $\eta_1=0.25$，$\eta_2=1.0$。由此可见，AGF-ADNet 将局部重构困难度、双机理分歧程度和全局分布漂移统一到同一异常评分框架下。检测重构分数反映正常模式建模难度，专家分歧分数反映不同恢复机理之间的一致性，分布漂移分数反映窗口统计属性是否脱离训练分布。三者结合后，异常评分不再依赖单一证据来源。

\subsection{时序平滑与阈值判别}

为抑制孤立波动对判别结果的干扰，对异常分数序列执行指数加权平滑：
\begin{equation}
    \tilde{s}_t = \alpha s_t + (1-\alpha)\tilde{s}_{t-1},
\end{equation}
其中 $\alpha \in (0,1]$ 为平滑系数，当前取值为 $0.3$。若测试数据由多个独立工况片段组成，则分别在各片段内部执行平滑，以避免跨片段信息泄漏。

在训练分数分布基础上，异常判别阈值定义为
\begin{equation}
    \theta
    =
    \max
    \left(
        \mu_{\mathrm{tr}} + \kappa \sigma_{\mathrm{tr}},
        \operatorname{Quantile}_{q}(\tilde{s}_{\mathrm{tr}})
    \right),
\end{equation}
其中，$\kappa=2.0$，$q=0.95$。据此得到初始预测结果
\begin{equation}
    \hat{y}_t =
    \begin{cases}
        1, & \tilde{s}_t \geq \theta, \\
        0, & \tilde{s}_t < \theta.
    \end{cases}
\end{equation}

考虑到工业异常通常具有一定持续时间，进一步对预测结果施加最短持续长度约束。若某一连续异常片段长度小于预设阈值 $L_{\min}$，则将该片段抑制为正常状态。在当前实验中，$L_{\min}=50$。这一后处理能够有效降低短时尖峰噪声引起的误报。

\section{本章小结}

本章围绕当前版本 AGF-ADNet 的结构与机理，对数据预处理、采样率元数据构造、双向预插补、自适应图学习、图重构专家、频域重构专家、采样率感知门控融合、Transformer 全局重构检测器以及异常评分机制进行了系统阐述。与传统先插补后检测的串联式方法相比，AGF-ADNet 通过联合优化实现了插补与检测的协同建模；同时，模型显式引入采样率类别和相位信息，使其能够更充分地适应多采样率工业过程中的结构性缺失、变量耦合以及周期动态特征。上述设计为后续实验验证和性能分析提供了统一的理论基础。
